% ============================================================
% Chapter 4: Proposed Method: Specification and Design
%   Section: The Fact-Checking Pipeline
%     Subsection: Evidence Retrieval
% Included from chapters/04-proposed-method/06-fact-checking-pipeline/fact-checking-pipeline.tex
% Source: docs/decisions/retrieval.md; retrieval/query_builder.py.
% ============================================================

\subsection{Evidence Retrieval}
\label{subsec:evidence-retrieval}

\paragraph{The failure the design prevents.} An early version built one
query per claim from its entities, figures and dates. A Spanish claim
that the Coyote's mail-order ACME purchases stood for post-war American
consumerism was searched as ``coyote ACME Estados Unidos Segunda
Guerra Mundial'', which every page explaining the Greek word
\emph{acme} matches; those pages were the best embedding match, the
model was shown three of them, and the claim came back \texttt{FALSE}
at 83\% confidence. Each stage had behaved as designed, and none could
produce the honest answer, which was \texttt{UNVERIFIED}.

\paragraph{Three queries, fused.} Each claim is now searched with a
planned set of three queries:

\begin{itemize}
    \item an \textbf{anchor} query --- entities, quoted figures and
    dates --- for precision on the subject;
    \item a \textbf{proposition} query --- the claim's own content
    words plus two anchors --- for what the claim asserts, which the
    anchor query lacks;
    \item a \textbf{refutation} query --- the anchors plus words that
    would contradict the claim, per language --- because every query
    built from a claim is phrased affirmatively and could otherwise only
    confirm it.
\end{itemize}

The three run concurrently and their ranked lists are fused by
reciprocal rank fusion \cite{cormack2009rrf}, which orders candidates by
how much the queries agree rather than by which query ran first, and
each candidate records which queries found it. When a claim lost its
subject on being separated from its article (``the programme trained
30 experts''), the headline words that the article's own keywords agree
on are restored to its queries, and the article's own URL is dropped
from the results, so an article cannot confirm itself.

\paragraph{Two sources of candidates.} The queries go to SearXNG, which
consults a closed list of engines measured one by one --- general web
engines and, at half weight, four academic indexes (arXiv, Crossref,
Semantic Scholar and PubMed) that keep answering when the web engines
rate-limit. In parallel, the system's own corpus of analysed articles is
searched by embedding, always excluding the article itself.

\paragraph{A funnel.} Reading a page in full is the most expensive
operation of the pipeline, so candidates are first pre-ranked cheaply
on their title and snippet, with the same semantic and lexical split
the ranking uses, and only the best eight are fetched through the
extraction cascade. Evidence pages never escalate to the headless
browser.

\paragraph{Search unavailable is not a verdict.} If every query of a
claim fails --- SearXNG unreachable, or answering empty while reporting
its engines down --- the claim is marked \texttt{searchUnavailable}
and stops here with a note that the web was never asked. Its verdict is
still \texttt{UNVERIFIED}, but it no longer looks like a judgement about
the claim, and the evaluation sets it apart. Before giving up, a query
SearXNG cannot answer is sent once to a second, keyless route.
